% Generated draft: requires mathematical and source review.
\documentclass[11pt]{article}
\usepackage[margin=25mm]{geometry}
\usepackage{fontspec}
\setmainfont{Latin Modern Roman}
\usepackage{amsmath,amssymb,mathtools,unicode-math}
\setmathfont{Latin Modern Math}
\usepackage{enumitem,needspace,xurl,hyperref}
\setlength{\parindent}{0pt}
\setlength{\parskip}{5pt}
\setlength{\emergencystretch}{3em}
\newcommand{\N}{\mathbb{N}}
\newcommand{\Z}{\mathbb{Z}}
\newcommand{\Q}{\mathbb{Q}}
\newcommand{\R}{\mathbb{R}}
\newcommand{\C}{\mathbb{C}}
\newcommand{\F}{\mathbb{F}}
\newcommand{\A}{\mathbb{A}}
\newcommand{\PP}{\mathbb{P}}
\newcommand{\E}{\mathbb{E}}
\newcommand{\eps}{\varepsilon}
\newcommand{\dd}{\,\mathrm{d}}
\begin{document}
\section*{UnsolvedMath ID: 30005757}
\textbf{Source problem code:} catalogue entry
\subsection{Definitions and mathematical statement}
A finite permutation of length \(n\) is a bijection from \(\{1,\ldots,n\}\) to itself, written in one-line notation. For finite permutations \(\sigma\) and \(\pi\), write \(\sigma\leq\pi\) if some subsequence of \(\pi\) has the same relative order as \(\sigma\). A permutation class is a set of finite permutations closed downward under this relation. For a set \(B\) of permutations, define
\[
\operatorname{Av}(B)=\{\pi: \text{there is no }\beta\in B\text{ with }\beta\leq\pi\}.
\]
For a single forbidden permutation, write \(\operatorname{Av}(\beta)=\operatorname{Av}(\{\beta\})\). The class of separable permutations is
\[
\mathcal S=\operatorname{Av}(\{2413,3142\}).
\]
Here a subclass of \(\mathcal S\) means a permutation class \(\mathcal C\subseteq\mathcal S\), not an arbitrary subset.

Let \(G\) be the group of permutation symmetries generated by reversal, complementation, and inversion. For \(\pi=\pi_1\cdots\pi_n\), reversal gives \(\pi_n\cdots\pi_1\), complementation gives \((n+1-\pi_1)\cdots(n+1-\pi_n)\), and inversion gives the inverse bijection. For \(g\in G\), put \(g(\mathcal C)=\{g(\pi):\pi\in\mathcal C\}\).

Write \(\mathcal C_n\) for the permutations of length \(n\) in \(\mathcal C\), and interpret its generating function as the ordinary length generating function
\[
F_{\mathcal C}(x)=\sum_{n\geq0}|\mathcal C_n|x^n\in\Z[[x]],
\]
with the empty permutation included in each nonempty class. A series is rational if it is the expansion at zero of \(P(x)/Q(x)\), where \(P,Q\in\Q[x]\) and \(Q(0)\neq0\).

The supplied Question 4 asks whether the following implication holds for every permutation class \(\mathcal C\):
\[
\left[\mathcal C\subseteq\mathcal S\ \text{ and }\ \exists g\in G\; g(\operatorname{Av}(231))\subseteq\mathcal C\right]
\quad\Longrightarrow\quad
F_{\mathcal C}(x)\notin\Q(x).
\]
Thus the containment hypothesis concerns an entire avoidance class or its entire symmetry image.

The supplied surrounding discussion concerns separate notions. A permutation class is well-quasi-ordered under pattern containment if every infinite sequence \(\pi_1,\pi_2,\ldots\) in it has indices \(i<j\) with \(\pi_i\leq\pi_j\); for finite permutations this is equivalent to having no infinite antichain. A class is strongly algebraic if every permutation subclass has a generating function algebraic over \(\Q(x)\). The supplied context asserts that a class which is not well-quasi-ordered cannot be strongly algebraic, and mentions analogous obstructions for the properties that every subclass has a D-finite or D-algebraic generating function. It labels the equivalence between strong algebraicity and well-quasi-ordering as a false conjecture attributed to Vatter, reference [9], Conjecture 12.3.4, and attributes its disproof to Brignall and Vatter, reference [6]. These contextual assertions are supplied source claims, not additional claims established by this draft. The phrase ``such classes'' in the contextual well-quasi-ordering assertion has an unclear antecedent, and the counterexample description is truncated after ``uncountably''; neither is used as a hypothesis of the displayed question.

The supplied dated assessment labels the record partially solved while stating that this exact rationality question remains unresolved. It describes a known theorem as proving the opposite implication, without supplying its precise statement, and distinguishes the disproved strong-algebraicity conjecture from the question above. These are supplied curation claims requiring independent review.
\subsection{Short English statement}
Must every permutation class contained in the separable permutations have a nonrational ordinary counting generating function whenever it contains all \(231\)-avoiding permutations, or a symmetry image of that entire class? The supplied assessment says this exact question remains unresolved; its partially solved label and the separate disproved strong-algebraicity conjecture do not constitute a solution.
\subsection{Sources}
\begin{itemize}
\item Local project registry: supplied complete source record for permanent UnsolvedMath ID \(30005757\), title \emph{Rationality of Separable Permutation Subclasses}, exact source code \texttt{catalogue entry}. Project registry paths are \texttt{data/unsolvedmath/problems.json} and the supplementary \texttt{catalog/problems.json}; neither file was independently inspected for this draft. The extracted report question is numbered Question 4. No individual proposer is identified in the supplied record.
\item Supplied report citation: \emph{Mini-Workshop: Permutation Patterns} (2024), Oberwolfach Reports, pp.\ 273--308, DOI 10.4171/owr/2024/6. Supplied source URL and literature-source label ``original OWR report'': \url{https://doi.org/10.4171/owr/2024/6}.
\item Supplied literature source labeled ``separable-subclass theorem'': \url{https://arxiv.org/abs/1007.1014}. The record does not reproduce the precise theorem or its hypotheses.
\item Supplied literature source labeled ``wqo counterexample paper'': \url{https://arxiv.org/abs/2211.12397}. The surrounding excerpt attributes a disproof of the separate strong-algebraicity conjecture to Brignall and Vatter, reference [6]. It attributes that conjecture to Vatter, reference [9], Conjecture 12.3.4; complete bibliographic details for reference [9] are not supplied.
\item Literature assessment date: \texttt{2026-08-22}. This is supplied dated curation, not independently verified here. It labels the record partially solved, says the exact rationality question remains unresolved, describes the known theorem as proving the opposite implication, and distinguishes a separate disproved strong-algebraicity conjecture. The record also supplies the label \texttt{corrected\_verified} and a claim that the official report was checked directly; this draft records that provenance claim without claiming a new source visit or verification.
\end{itemize}
\end{document}
